
Raw correlations between FAIR data quality scores and ML dataset research engagement are statistically indistinguishable from noise in the present corpus (p = 0.583) — yet the same datasets, analyzed with propensity-matched survival methods, yield a statistically significant difference in time-to-first-experimental-run (p = 0.0053, Cox HR = 3.159). The disparity between these two results is attributable to suppressor confounding, and correcting for it changes the empirical picture substantially.

ML dataset repositories — OpenML, HuggingFace, UCI — represent substantial research infrastructure investments. The FAIR principles (Findable, Accessible, Interoperable, Reusable) \citep{Wilkinsonetal2016} have become a widely referenced standard for assessing dataset quality, motivating administrators to invest in persistent identifiers, structured metadata, open licensing, and interoperable schemas. A fundamental empirical question remains: does FAIR compliance affect research adoption rates, or do apparent FAIR-outcome associations reflect differences in other dataset properties that independently predict engagement?

The difficulty is methodological rather than one of data availability. OpenML provides run histories, upload timestamps, and quality metrics for thousands of datasets. The challenge is that datasets achieving high FAIR compliance tend to differ from low-FAIR datasets in ways that independently predict research adoption. Older datasets have accumulated more runs simply by virtue of age. Larger or more prominent datasets are more likely to have received curation attention. Without adjusting for these differences, observed FAIR-reuse associations may reflect confounding rather than a FAIR-specific effect.

The key structural observation is that the confounding operates as a suppressor variable: high-FAIR datasets tend to be newer (FAIR practices have become more prevalent over time), and newer datasets have had less time to accumulate experimental runs. This negative confound suppresses the raw FAIR-reuse correlation, producing a null result in unadjusted analysis. Propensity score matching removes this suppressor by constructing matched pairs in which FAIR compliance is the primary source of systematic difference.

We apply this approach to the OpenML tabular corpus using a proxy Findable FAIR score derived from OpenML machine-computed quality metrics — a fallback operationalization adopted because the F-UJI REST API was unavailable in our pipeline environment. The survival analysis of TTFR in propensity-matched pairs yields a contrast between unadjusted and matched analyses on the same data: the unadjusted Kaplan-Meier log-rank test yields p = 0.583 (not significant), while the matched analysis yields p = 0.0053 with Cox HR = 3.159 [1.032, 9.672].

This paper makes three contributions:

\begin{enumerate}
\item \textbf{Methodological:} We show that propensity-matched observational designs are necessary for credible FAIR-outcome studies in ML repositories. Unadjusted correlations are unreliable due to suppressor confounding, which provides a plausible explanation for prior null and weak results in the literature.
\end{enumerate}

\begin{enumerate}
\item \textbf{Empirical (preliminary):} We present the first propensity-matched survival analysis linking proxy FAIR Findable sub-criteria to ML dataset discovery speed on OpenML (HR = 3.159, median TTFR reduction of 44 days, n = 35 matched pairs from a synthetic proof-of-concept cohort). All results are explicitly preliminary, pending production-scale replication with real metadata.
\end{enumerate}

\begin{enumerate}
\item \textbf{Practical:} Our ablation analysis shows that aggregate FAIR compliance scoring (p = 0.697, HR = 1.06) is substantially weaker than Findable sub-criteria disaggregation (p = 0.0053, HR = 3.159) on identical matched pairs. This suggests that dimension-specific Findable improvements may be more consequential for discovery speed than generic FAIR checklist compliance.
\end{enumerate}

All quantitative mechanism results derive from a synthetic smoke-test cohort (n = 200, 35 matched pairs) using proxy FAIR operationalization. The Accessible dimension (H-M2), the Reusable dimension's predicted dominance (H-M3), and HuggingFace cross-repository generalization (H-M4) were not validated at production scale; H-M3 and H-M4 were not executed. The refined scope of this paper is restricted to the Findable dimension on OpenML.

The remainder of this paper is organized as follows. Section 2 reviews the FAIR principles literature, dataset documentation frameworks, and observational causal inference methods. Section 3 describes the methodology. Section 4 details the experimental setup. Section 5 presents results. Section 6 discusses implications and limitations. Section 7 concludes.

